\documentclass[11pt]{article}
\usepackage{amsmath}
\usepackage{booktabs}
\usepackage[margin=1in]{geometry}

\title{Bayesian Additive Player and Season Models}
\author{}
\date{}

\begin{document}
\maketitle

\section{Observed Data}

The modelling unit is an observed player-season row. Missing player-seasons are not imputed. Let
$y_n \in \mathbb{R}^P$ be the scaled feature vector for row $n$, with player index $i[n]$
and season index $j[n]$.

\section{Real-Data Diagonal Additive Model}

\begin{align}
y_n &= \mu + a_{i[n]} + b_{j[n]} + \epsilon_n, \\
a_i &\sim \mathcal{N}_P(0, \Sigma_a), \\
b_j &\sim \mathcal{N}_P(0, \Sigma_b), \\
\epsilon_n &\sim \mathcal{N}_P(0, \Sigma_\epsilon).
\end{align}

For the real-data model, all covariance matrices are diagonal:

\begin{align}
\Sigma_a &= \operatorname{diag}(\sigma_{a1}^2,\ldots,\sigma_{aP}^2), \\
\Sigma_b &= \operatorname{diag}(\sigma_{b1}^2,\ldots,\sigma_{bP}^2), \\
\Sigma_\epsilon &= \operatorname{diag}(\sigma_{\epsilon1}^2,\ldots,\sigma_{\epsilon P}^2).
\end{align}

The Stan implementation uses non-centered standard-normal player and season effects. Columns are centered in transformed parameters:

\begin{align}
a_{ip} &= \sigma_{ap}\left(z^a_{ip} - \bar{z}^a_{\cdot p}\right), \\
b_{jp} &= \sigma_{bp}\left(z^b_{jp} - \bar{z}^b_{\cdot p}\right).
\end{align}

This leaves $\mu_p$ as the population-level feature mean in scaled units.

\section{Simulation Design}

Simulated data use the observed player-season rows and index structure but generate player and
season effects from known low-rank plus diagonal covariance:

\begin{align}
\Sigma_a &= \Lambda_a \Lambda_a^\top + \operatorname{diag}(\Psi_a), \\
\Sigma_b &= \Lambda_b \Lambda_b^\top + \operatorname{diag}(\Psi_b), \\
\Sigma_\epsilon &= \operatorname{diag}(\sigma_\epsilon^2).
\end{align}

The default ranks are $r_a = 2$ and $r_b = 1$, with $r_a > r_b$. Random effects are generated as

\begin{align}
a_i &= \Lambda_a \eta_i^a + \operatorname{diag}(\sqrt{\Psi_a}) u_i^a, \\
b_j &= \Lambda_b \eta_j^b + \operatorname{diag}(\sqrt{\Psi_b}) u_j^b,
\end{align}

where the latent scores and uniqueness terms are standard normal and column-centered.

\section{Low-Rank Recovery Model}

The recovery model matches the simulation:

\begin{align}
y_n &= \mu + a_{i[n]} + b_{j[n]} + \epsilon_n, \\
a_i &= \Lambda_a \eta_i^a + \operatorname{diag}(\psi_a) u_i^a, \\
b_j &= \Lambda_b \eta_j^b + \operatorname{diag}(\psi_b) u_j^b.
\end{align}

Loadings use pure-anchor constraints. For a rank-$Q$ loading matrix, the first $Q$ variables are
positive own-factor anchors and their off-factor loadings are fixed to zero. Non-anchor rows are
estimated freely. This removes sign and rotation ambiguity enough for direct comparison with
the saved simulation truth.

\section{Recovery Targets}

The diagnostics compare posterior summaries against saved values of
$\Lambda_a$, $\Lambda_b$, $\Psi_a$, $\Psi_b$, $\operatorname{diag}(\Sigma_a)$,
$\operatorname{diag}(\Sigma_b)$, $\sigma_\epsilon$, and $\mu$.

\end{document}
